\documentclass[11pt]{article}
\usepackage[T1]{fontenc}
\usepackage[utf8]{inputenc}
\usepackage[english]{babel}
\usepackage[margin=1in]{geometry}
\usepackage{amsmath,amssymb}
\usepackage{enumitem}
\usepackage{microtype}

\ifdefined\pdfcatalog
  \pdfcatalog{/Lang (en-US)}
\fi

\setlength{\parindent}{0pt}
\setlength{\parskip}{0.55em}
\setlength{\emergencystretch}{2em}
\setlist{nosep,leftmargin=*}

\newcommand{\findinglabel}[1]{\textbf{#1}}
\newcommand{\classification}[1]{\textit{Classification: #1.}}

\title{Technical Review of ``Early-Stopping Activation Steering for Factual Preference Alignment in Vietnamese Medical Language Models''}
\author{Manuscript review}
\date{August 28, 2026}

\begin{document}
\selectlanguage{english}
\maketitle

\section*{Overall assessment}

The current revision fixes several concrete defects.  It now provides the bounded-generation $2\times2$ table, makes the matched-dose definition sign-consistent, explains RefPref tie handling and its proxy status, uses a coherent 100-direction placebo summary, and compiles successfully with the natural-completion table in a full-width float.  The BM25 comparison is also substantially more informative because it includes retrieval parameters, an oracle condition, paired discordant counts, and latency dispersion.

The manuscript nevertheless still requires major revision.  The natural-completion significance results cannot be reproduced because the discordant counts, test variants, bootstrap specification, and Holm ordering are absent.  The placebo experiment supports directional specificity, but the Results still juxtaposes a descriptive $5.15\sigma$ distance with a finite-resolution empirical $p=0.0099$ in a way that can be mistaken for Gaussian five-sigma evidence.  Activation norms do not establish saturation prevention, natural completion shows a multi-metric trade-off rather than consistent improvement by Linear Decay, and RefPref remains inadequate for clinical correctness or hallucination claims.  Important dataset, selection, metric, hook, timing, and systems details also remain missing.

\section*{Major technical findings}

\subsection*{1. Natural-completion inference is still not reproducible}

\classification{Definite statistical-reporting omission and likely test inconsistency}

\findinglabel{Location.} Natural Termination Validation; Table~\texttt{tab:mcnemar}; Discussion; Conclusion.

\findinglabel{Problem.} The new bounded-generation table is internally reconstructable: $b=16$, $c=37$, a $4.20$-point paired gain, and exact McNemar $p\approx0.0055$.  The natural-completion table, however, reports only marginal totals, confidence intervals, and $p$-values.  It does not give the four paired cells or even $b$ and $c$, so none of its tests can be independently verified.

The displayed natural values also appear to use different possible test conventions.  The Hard-Cutoff effect of 26 records and $p\approx0.00086$ is compatible with $b=16,c=42$ under an exact test.  The Continuous value $p=0.0412$ is compatible with a continuity-corrected asymptotic calculation for one possible transition table, whereas the Linear-Decay value $p=0.1250$ does not identify a unique exact or asymptotic calculation from the marginal totals.  The manuscript does not state the test variant.  It likewise omits the bootstrap unit, number of replicates, seed, interval method, and the ordered raw values used for Holm adjustment.

\findinglabel{Why it matters.} The natural-completion significance claims are central to schedule selection, but marginal accuracies do not determine paired McNemar statistics.  Mixing or leaving unspecified test variants can change whether a comparison crosses the significance threshold.

\findinglabel{Suggested fix.} Report all four paired cells for every natural comparison, identify $b$ and $c$, and state the exact software calls and exact versus asymptotic convention.  Document the bootstrap unit, replicates, seed, and interval construction, and show the raw and Holm-adjusted $p$-values in adjustment order.  Generate tables and prose from one versioned analysis artifact.

\subsection*{2. The placebo result should not be presented as Gaussian five-sigma evidence}

\classification{Unsupported strength of statistical interpretation}

\findinglabel{Location.} Abstract; contribution~4; placebo Methods and Results; Table~\texttt{tab:baselines}; Conclusion.

\findinglabel{Problem.} The revised arithmetic is coherent:
\[
\frac{77.20-55.82}{4.15}\approx5.15,
\qquad
\frac{1}{100+1}\approx0.0099.
\]
The first number is a descriptive standardized distance using the observed dispersion of 100 bounded accuracy values; it is not automatically a calibrated Gaussian significance level.  Accuracy is bounded, distributional normality is not established, and the empirical randomization result has minimum plus-one resolution $1/101$.  The Abstract, contribution list, and Conclusion now call the distance descriptive, but the Results still reports ``$+5.15\sigma$'' beside $p=0.0099$ without the same qualification.  The meaning of $4.15\%$ should also be stated explicitly at the table, and the single placebo BERTScore value 0.6945 is not identified as a mean, pooled score, or one direction's result.

\findinglabel{Why it matters.} Direction specificity is useful evidence, but $5.15\sigma$ and an empirical $p=1/101$ are not equivalent inferential statements.  The current presentation can materially exaggerate certainty.

\findinglabel{Suggested fix.} State ``0 of 100 sampled directions exceeded the target; plus-one empirical $p=1/101$'' and call 5.15 only a descriptive standardized distance everywhere.  Define each $\pm$ quantity and release direction-level RefPref and BERTScore values with seeds.  Use more random directions if finer randomization resolution is required.

\subsection*{3. Activation-norm dynamics do not establish saturation prevention}

\classification{Unsupported mechanism claim and implementation ambiguity}

\findinglabel{Location.} Abstract; Introduction; contribution~3; Early-Stopping Steering Hook \& Saturation Diagnostics; Empirical Activation Magnitude Dynamics.

\findinglabel{Problem.} The Results now correctly calls 51.20 a $-3.01$-unit undershoot and cautions that scalar norms do not describe manifold geometry.  The heading and higher-level claims nevertheless continue to frame the analysis as saturation prevention.  At $t=50$, Linear Decay is farther from the displayed 54.21 baseline than Continuous Steering is:
\[
|51.20-54.21|=3.01,
\qquad
|56.18-54.21|=1.97.
\]
No prespecified acceptable band, clipping measure, responsiveness test, directional-collapse criterion, or behavioral stability definition makes the larger undershoot evidence of stabilization.

Only selected time points are reported, without dispersion, confidence bands, or per-step survivor counts.  After $t>K$, the early-stop schedules contain no direct additive perturbation, but free-generation histories differ between conditions; later state differences therefore conflate intervention effects with different preceding tokens.  The Introduction's claims about loops and degraded syntax are not connected to a controlled diagnostic or reported error analysis.

\findinglabel{Why it matters.} The central novelty is justified as preventing an internal failure, whereas the current experiment establishes only post-hook norm elevation and undershoot on divergent generation contexts.

\findinglabel{Suggested fix.} Define saturation operationally and test it on fixed-token replay or teacher-forced contexts.  Report complete pre/post-hook trajectories, uncertainty, survivor counts, projection onto $v_{\text{steer}}$, cosine drift, distributional distance, and response to an additional controlled perturbation.  Link these measurements to repetition and factual errors.  Until then, use ``activation-norm dynamics'' rather than ``saturation prevention'' or ``stabilization.''

\subsection*{4. Natural completion shows a trade-off and does not select Linear Decay as the best schedule}

\classification{Unsupported comparative conclusion and remaining contradiction}

\findinglabel{Location.} natural-completion Results; Table~\texttt{tab:long\_gen}; Discussion; Conclusion.

\findinglabel{Problem.} Every steering schedule raises RefPref but lowers reference BERTScore and raises Rep-4 relative to baseline.  Hard Cutoff has the highest RefPref (70.60\%), baseline has the highest BERTScore (0.6779) and lowest Rep-4 (4.10\%), and Linear Decay has the highest repetition (5.37\%).  The Hard-Cutoff Rep-4 increase to 5.35\% is approximately 30.5\% relative, so calling it ``minor'' requires uncertainty or a practical threshold.  The natural table labels Linear Decay versus baseline nonsignificant at $p=0.1250$, while the Conclusion still opens by saying that linear decay yields consistent improvements.

The opening Results overview also continues to state that both baseline and early-stopping generations reach EOS in 100.00\% of cases, although Continuous Steering and Linear Decay are reported at 99.80\% elsewhere.

\findinglabel{Why it matters.} The data neither demonstrate uniform output-quality improvement nor establish the proposed decay schedule as the preferred natural-generation intervention.  Post-hoc selection among schedules and metrics can overstate the evidence.

\findinglabel{Suggested fix.} Prespecify one primary endpoint and schedule on validation data, then evaluate that frozen choice once on the test set.  Report paired uncertainty for RefPref, BERTScore, and Rep-4, with noninferiority margins if degradation is claimed to be acceptable.  Use metric-specific language, the 99.80\%--100.00\% EOS range, and a neutral account of Hard Cutoff versus Linear Decay.

\subsection*{5. RefPref remains insufficient for clinical factual-correctness claims}

\classification{Unsupported interpretation of a disclosed proxy}

\findinglabel{Location.} title; Abstract; Introduction; contributions; Results; Discussion; Conclusion; keywords.

\findinglabel{Problem.} The Methods now appropriately states that RefPref is a semantic proxy and defines ties as zero while retaining them in the denominator.  Elsewhere, however, the vector is called ``truthfulness,'' the endpoint is described as accuracy or factual alignment, and the manuscript retains hallucination-mitigation and clinical-decision-support framing.  Relative BERTScore similarity to $y^+$ rather than one synthetic $y^-$ cannot establish whether an answer contains a wrong dose, unit, negation, contraindication, or interaction mechanism.  The natural results also show RefPref increasing while similarity to the reference decreases.

\findinglabel{Why it matters.} Clinical safety, harmful-error reduction, and deployment suitability require direct correctness and harm labels.  Semantic preference is not interchangeable with those endpoints.

\findinglabel{Suggested fix.} Add blinded clinician-adjudicated correctness and harmful-error rates, category stratification, reviewer qualifications, agreement, and adjudication.  Validate RefPref against those labels.  Otherwise narrow the title, Abstract, keywords, and conclusions to reference preference and remove clinical-accuracy, truthfulness, and hallucination claims.

\subsection*{6. The stronger RAG experiment still supports only one BM25 configuration}

\classification{Likely weak baseline and incomplete systems interpretation}

\findinglabel{Location.} Experimental Setup; RAG Results; Table~\texttt{tab:baselines}; Discussion; Conclusion.

\findinglabel{Problem.} The added oracle condition, BM25 parameters, paired counts, and latency dispersion are useful.  The paired RAG comparison is internally consistent: $b=21$, $c=64$ gives an 8.60-point net difference and a very small exact McNemar value.  However, Recall@1 is only 19.20\%, while Oracle Gold-Context RAG reaches 89.40\%, above the 77.20\% steering result.  This pattern implicates the chosen retriever rather than RAG as a method class.  Passage construction, query formation, relevance annotation, retrieval-failure handling, and stronger lexical or dense Vietnamese baselines remain absent.

The latency $\pm$ term is now defined as an across-query sample standard deviation, but one observation per query does not separate run-to-run variance from query heterogeneity.  The claimed ``zero-context-memory advantage'' is not accompanied by peak-memory or KV-cache measurements.  Identical displayed steering and baseline timing values also do not establish equivalence without repeated component-wise measurements and a margin.

\findinglabel{Why it matters.} The paper can compare steering with this BM25 pipeline, but it cannot yet generalize to RAG or support unmeasured memory and equivalence claims.

\findinglabel{Suggested fix.} Describe the conclusion as specific to the tested BM25 setup.  Publish passage construction, queries, relevance labels, and failure policy; add stronger retrieval baselines and retrieval-conditioned generation results.  Measure retrieval, prefill, decoding, total latency, peak accelerator memory, and KV-cache allocation over repeated runs, with an equivalence margin for ``negligible'' overhead.

\subsection*{7. Dataset construction and model-selection controls remain incomplete}

\classification{Reproducibility limitation and likely leakage risk}

\findinglabel{Location.} benchmark construction; vector estimation; layer and coefficient selection; Experimental Setup.

\findinglabel{Problem.} ``Expert-structured'' remains undefined.  The manuscript gives no annotation protocol, counter-answer construction rules, clinician validation, deduplication method, source identifiers, licensing statement, or source/template-grouped split manifest.  A question-disjoint split does not by itself prevent leakage through shared drugs, formulary passages, or templates.  The selection rule and seed for the 500-item evaluation subset are absent, and the relationship between the three reported category counts and the full 14,700-item corpus is unclear.

Layer~8, $\alpha_0=18$, and $K=16$ are used as primary settings without a complete validation-selection procedure.  The layer sweep and ablations are reported on the same benchmark used for the headline result, creating a risk of test-set selection.  Small BERTScore differences among schedules and dose controls have no paired intervals, multiplicity handling, or practical-equivalence threshold.

\findinglabel{Why it matters.} Leakage and post-hoc hyperparameter choice can inflate performance, while missing provenance prevents independent reconstruction of the benchmark.

\findinglabel{Suggested fix.} Add a dataset card, item/source identifiers, licenses, annotation and adjudication procedures, grouped split manifests, evaluation sampling code and seed, and category counts for every split.  Freeze layer, coefficient, and schedule using a validation set before one test evaluation.  Report paired uncertainty and exploratory multiplicity handling for ablations.

\subsection*{8. Metric, vector, hook, and generation details remain implementation-critical omissions}

\classification{Reproducibility limitation and implementation ambiguity}

\findinglabel{Location.} vector construction; metric definitions; steering hook; Experimental Setup; placebo table.

\findinglabel{Problem.} The direction is extracted from the final token of a completed answer and injected into early autoregressive states without prompt-only, initial-token, pooled-position, or independent-error-template controls.  The exact hook module, zero- versus one-based layer convention, tensor positions, prompt-forward treatment, KV-cache behavior, and pre/post-hook logging order are not specified.  Rep-4 lacks a complete formula and tokenizer definition; Vietnamese ROUGE tokenization, BERTScore model/configuration and rescaling, and statistical software calls are incomplete.  The placebo BERTScore value 0.6945 has no aggregation definition or across-direction dispersion.

\findinglabel{Why it matters.} These choices can materially change the learned direction, intervention, metrics, and statistical results even when package versions are fixed.

\findinglabel{Suggested fix.} Release executable hook code or pseudocode, exact module/tensor/index conventions, prompt and cache handling, all generation settings and seeds, complete metric definitions, and statistical calls.  Report direction-level placebo metrics and test alternative vector extraction positions and independently constructed negative answers.

\subsection*{9. Layer, ablation, and deployment interpretations exceed the reported evidence}

\classification{Unsupported explanatory and deployment claims}

\findinglabel{Location.} layer sweep; ablation study; Discussion; Conclusion.

\findinglabel{Problem.} The layer sweep reports only a narrow aggregate range without per-layer paired uncertainty, yet it is used to infer distributed truthfulness or robust subspace structure.  Ablation BERTScore differences are very small and lack confidence intervals or correction for multiple comparisons.  The Discussion attributes behavior to attention/template mechanisms and suggests that repetition penalties will preserve the RefPref gain without testing either claim.  Privacy, local-hospital deployment, and small-language-model language are not supported by a threat model, measured resource envelope, or explicit model-size definition.

\findinglabel{Why it matters.} Exploratory patterns and proposed mitigations are being presented as explanatory or deployment evidence.

\findinglabel{Suggested fix.} Label these observations as hypotheses, add paired layer/ablation inference on held-out data, and experimentally test proposed repetition controls.  Provide a deployment threat model, hardware/memory measurements, and an explicit definition of ``small'' before making privacy or local-deployment claims.

\section*{Equation, notation, sign, branch, and dimensional audit}

The direction $\mu_l^+-\mu_l^-$ and intervention $+\alpha(t)v_{\text{steer}}^{(l)}$ are sign-consistent with the intended positive intervention.  Unit normalization makes the direction dimensionally compatible with the hidden state when $\alpha(t)$ is an activation-scale coefficient.  The linear schedule and the definitions of $D_{\text{continuous}}$, $D_{\text{cutoff}}$, and
\[
D_{\text{decay}}=\frac{K+1}{2}|\alpha_0|
\]
are arithmetically correct.  The revised matched-dose expression using $\operatorname{sgn}(\alpha_0)$ correctly matches the absolute dose, and the displayed positive coefficients 0.6375, 0.7650, and 0.8500 are correct for $K=16$ and $T=200$.  No inverse-trigonometric expressions, coordinate transformations, or branch/quadrant choices occur.

Remaining notation issues are implementation-facing rather than algebraic.  The symbol $h_l(x_i,y_i)$ is first defined as the final-token representation of a completed prompt--answer sequence, while $h_l^{(t)}$ later denotes a generation-time residual stream; these roles should be distinguished explicitly.  Pre-hook $h_l^{(t)}$ and post-hook $\hat h_l^{(t)}$ should remain separate everywhere.  $N$ is overloaded for training examples, category sizes, evaluation prompts, and random directions.  Decoding steps are discrete and should be written $t\in\{1,\ldots,100\}$ rather than $t\in[1,100]$.

\section*{Meaningful editorial, compilation, and reference findings}

\subsection*{The source compiles, but minor table overflow and widespread loose typesetting remain}

\findinglabel{Location.} Table~\texttt{tab:mcnemar\_bounded}; clean two-pass pdfLaTeX log.

\findinglabel{Problem.} The previous natural-table float error and severe one-column overflow are fixed.  A clean two-pass build succeeds.  The bounded McNemar table is still approximately 1.77~pt too wide, and the log reports numerous high-badness underfull boxes in the Abstract, contribution list, environment description, and dense technical paragraphs.

\findinglabel{Why it matters.} The remaining overflow is small but avoidable, while repeated loose lines reduce the visual quality of a compact conference manuscript.

\findinglabel{Suggested fix.} Shorten or wrap the bounded-table headings, reduce intercolumn padding without harming readability, and revise long unbreakable technical phrases.  Rebuild twice and inspect the final PDF and log.

\subsection*{Citation fit remains uneven}

\findinglabel{Location.} Introduction; bibliography entries \texttt{ref10}, \texttt{ref11}, and \texttt{ref12}.

\findinglabel{Problem.} \texttt{ref10} is the LoRA paper rather than direct evidence that SFT mitigates the stated medical hallucinations.  \texttt{ref11} studies irrelevant retrieved context but does not establish that internal activation pathways cause evidence ignoring.  \texttt{ref12} concerns catastrophic forgetting generally rather than the asserted medical-language-model failure mode.

\findinglabel{Why it matters.} Direct evidence, general motivation, and mechanistic hypotheses are presented with similar certainty.

\findinglabel{Suggested fix.} Add directly matched sources or qualify the corresponding statements as motivation and hypotheses.

\section*{Proofing sweep}

The bundled mechanical scan was run on both the source and the clean compiled PDF.  Its only hit was the semicolon before ``Recall@3,'' which is not a capitalization defect.  Manual source, equation-adjacent, table, and bibliography checks identified these bounded corrections:

\begin{itemize}
    \item \textbf{Results overview:} replace the universal 100.00\% EOS statement with the condition-specific 99.80\%--100.00\% range.
    \item \textbf{Placebo Results:} add ``descriptive standardized distance'' beside $5.15\sigma$ and define the table's BERTScore aggregation.
    \item \textbf{Activation terminology:} replace the remaining ``Saturation Diagnostics'' heading and write ``activation $\ell_2$ norm'' rather than ``Activation L2 norm.''
    \item \textbf{Timing units:} use nonbreaking spaces in values such as 7.13~s and define $\pm$ quantities at their first table occurrence.
    \item \textbf{Bounded McNemar table:} remove the remaining approximately 1.77~pt overflow by shortening or wrapping headings.
\end{itemize}

The bibliography was checked for duplicated punctuation, malformed quotation punctuation, inconsistent capitalization, and obvious citation-format glitches; no additional mechanical defect was identified.

\section*{Items requiring external verification}

\begin{itemize}
    \item Verify all formulary-derived reference answers and synthetic counter-answers against the official source, current clinical guidance, qualified clinician review, and applicable reuse permissions.
    \item Validate multilingual BERTScore for Vietnamese clinical negation, quantities, units, contraindications, and interaction mechanisms before interpreting RefPref as factual alignment.
    \item Verify novelty against current token-dependent, scheduled, and dose-controlled activation-steering research.
    \item Verify the characterization of Qwen2.5-7B-Instruct as an SLM and the privacy/local-hospital claims using explicit definitions, a threat model, and measured resource requirements.
    \item Verify BM25 and stronger Vietnamese retrieval baselines before generalizing the RAG comparison.
\end{itemize}

\section*{Highest-priority revision sequence}

\begin{enumerate}
    \item Publish every natural-completion paired cell and regenerate the tests, confidence intervals, and Holm adjustment under one documented procedure.
    \item Align the activation claims with the limited norm diagnostic, or add controlled multidimensional mechanism experiments.
    \item Add clinician-adjudicated correctness and harm endpoints, or narrow all clinical, truthfulness, accuracy, and hallucination claims to RefPref.
    \item Present natural completion as a metric trade-off and resolve the primary Linear-Decay versus Hard-Cutoff choice using validation data.
    \item Reframe the placebo result as finite-resolution directional evidence rather than Gaussian five-sigma significance.
    \item Delimit and strengthen the RAG comparison and add component-wise repeated timing and memory measurements.
    \item Complete dataset provenance, leakage controls, selection rules, metric/hook details, ablation uncertainty, and remaining table/layout fixes.
\end{enumerate}

\end{document}